% Root and B full literal audits passed; D independently verified both H3 controls.
\subsection{Connected six-state graphs opposite the right-zero two-band}

\begin{lemma}[Polynomial graph actions and a degree-two square]
\label{C:connected-graph-polynomial-square}
Let a finite monoid act on a connected graph, and let $X$ be
the cofiber of its action-bar projection. Suppose an actual
equivariant self-map induces an operator $F$ on
$W=H_1(\Gamma;\mathbb Q)$ with $(F-1)(F-2)=0$.
If its pullback multiplies a degree-two class $v$ by two,
then $v^2=0$.
\end{lemma}
\begin{proof}
The natural connected graph comparison identifies every
positive homology group of $X$ with a shifted derived tensor
of $W$. Functoriality and additivity make
$(f^*-1)(f^*-2)=0$ on every positive rational cohomology
group. The square is multiplied by four, so that polynomial
applied to it gives $(4-1)(4-2)v^2=0$.
\end{proof}

\begin{theorem}[The connected six-state right-zero obstruction]
\label{C:right-zero-six-connected}
Let $R=\{1,p,q\}$ be the identity-adjoined right-zero
two-band, and let $|S|\leq6$ in an $(R,M)$-biset.
If its natural left-action graph is connected, the collage
cannot have type $\mathbb CP^2$. There is no bound on $|M|$.
\end{theorem}
\begin{proof}
A singleton core gives a common fixed state. Connected
graphs with at most five states were treated in
Theorem~\ref{C:right-zero-five-connected}.
The remaining graph has either three vertices and three
edges, which is unicyclic and excluded by
Lemma~\ref{C:unicyclic-graph-correspondence}, or two vertices
and four edges. A three-vertex tree is contractible and
cannot occur either. Label the two core vertices $0,1$.
We assume the second rational cohomology is one-dimensional,
and write $\lambda$ for its nonzero invariant functional
on $W=H_1(\Gamma;\mathbb Q)$.
All graph quotients used below are actual equivariant maps.
When their first-homology coinvariants have the same
nonzero one-dimensional functional, the natural graph-Tor
comparison makes their cofiber pullback an isomorphism on
second rational cohomology.

\emph{No loops.}
A constant vertex map kills all first homology, so none
can occur. Without a vertex swap a core vertex is fixed.
Thus both cross-edge orientations occur and a swap occurs.
For three edges in one direction and one in the other,
any swap kills all differences among the three parallel
edges. Identifying them gives the two-edge circle and an
isomorphism on first-homology coinvariants. Its degree-three
map kills the square.

For two edges in each direction let $P$ be the two-dimensional
space of differences within the two parallel pairs. If an
actor folds one pair, then either a swap transfers that
relation to the other pair, or the square of a nonbijective
swap kills both differences. Thus $P$ has zero coinvariants.
If every actor is a permutation, but a nonidentity
vertex-fixing permutation occurs, it reverses one or both
parallel differences; a vertex swap again forces both to
vanish rationally. These cases reduce to the same two-edge
circle quotient. Otherwise the entire graph action image
is just an involution interchanging the vertices. It acts
freely on vertices and edges, and its quotient has one
vertex and two edges. Its first-homology coinvariants have
dimension two, contrary to the assumed second homology.

\emph{Loops at only one vertex.}
There must again be a vertex swap, or a core vertex is fixed
unless a constant actor at the other vertex occurs; that
actor would kill all first homology. The swap kills the
whole loop subspace, because the other vertex has no loops.
Any constant actor then has image in a subspace already
zero in coinvariants, so makes all coinvariants zero.
Collapse the loops. This is an equivariant graph quotient,
with at most three cross edges, and is an isomorphism on
first-homology coinvariants. The connected five-state
rational square obstruction, proved in the connected part
of Lemma~\ref{C:nine-five-state-graph-square}, applies.
That connected argument places no bound on $M$.

There are only three remaining loop distributions.

\emph{One loop at each vertex and two parallel cross edges.}
Write the loops as $x_0,x_1$ and the parallel edges as
$a,b:0\to1$. There is no vertex swap, so both constant
vertex maps must occur. Invariance forces
$\lambda(x_0)=\lambda(x_1)=c\ne0$ and every constant
actor sends both loops to the loop at its target.
Every core-fixing actor fixes the loops.
On $y=a-b$, invariance gives
$\lambda(y)/c\in\{-1,0,1\}$.

If that value is zero, identifying $a=b$ is equivariant
and preserves the nonzero coinvariant functional, hence
reduces to a connected five-state graph. If it is one,
every core-fixing actor is the identity and every constant
actor sends $a$ to the target loop and $b$ to the target
vertex. Define $f(b)=b$ and $f(a)=ab^{-1}a$.
Map each loop twice around itself with a constant pause
between the two traversals. Parameterize these as three
equal intervals, so the middle collapsed $b^{-1}$ segment
matches the pause exactly under every constant actor.
This is a strictly equivariant map acting as twice the
identity on $W$. The scalar graph lemma kills the square.
The value minus one is the same argument with $a,b$
interchanged.

\emph{One loop at each vertex and opposite cross edges.}
Write $a:0\to1$, $b:1\to0$, and $z=a+b$.
If there is no constant vertex actor, a nonzero class on
the loop subspace either vanishes under a collapsed loop
and its swapped mate, or all actors act through the single
simultaneous swap of loops and cross edges. The latter
case has two-dimensional first-homology coinvariants.
Without a swap there is a fixed core vertex. Thus a rank-one
square in the no-constant case pulls back from the two-edge
circle quotient and is zero.

With a constant actor, avoiding a fixed core vertex forces
both constant vertex types, possibly by composing with a
swap. Invariance now gives
$\lambda(x_0)=\lambda(x_1)=c\ne0$, and the loops are
never collapsed. Core-fixing actors are identities; a swap,
if present, is the unique simultaneous swap.
For a constant actor write $\delta_a,\delta_b\in\{0,1\}$
for whether the corresponding cross edge enters the target
loop rather than the target vertex. The sum
\[
 d=\delta_a+\delta_b=\lambda(z)/c
\]
is the same for every constant actor.

For $d=0$, square both loops and fix $a,b$. The resulting
map is strictly equivariant, its operator on $W$ has
eigenvalues one and two, and $\lambda$ is multiplied by
two. For $d=2$, square the loops and put
$f(a)=x_0a$, $f(b)=x_1b$, with equal halves on each
concatenation. Constant actors send both halves to the
same loop, and the swap preserves the construction.
Here $F(x_i)=2x_i$ and $F(z)=z+x_0+x_1$, so again
$(F-1)(F-2)=0$ and $\lambda F=2\lambda$.
Lemma~\ref{C:connected-graph-polynomial-square} handles
both cases.

If $d=1$ and only the pattern $(1,0)$ occurs, there is
no swap. Square the loops, put $f(a)=x_0a$, and fix $b$.
Then $F(z)=z+x_0$, giving the same polynomial and eigenvalue.
The pattern $(0,1)$ is symmetric.
Finally, if both patterns occur, closure forces all four
constant actors. The actual right-action image is the
five-element monoid consisting of these actors and the
identity, or the six-element monoid also containing the
simultaneous swap.

The corresponding fourteen- and fifteen-arrow literal
collages are stored in
\path{code/outputs/a_rightzero_six_opposite_edges.json}.
Each has a nonzero third homology class over $\mathbb F_2$,
certified by an explicit normalized three-cycle and a
three-cocycle pairing to one. Their normalized four-simplex
counts are respectively $992$ and $1859$.
The standalone checker
\path{code/d_verify_rightzero_six_H3.py} reconstructs both
literal categories from their corner tables and actions,
checks all typed associativity triples, and independently
enumerates every normalized chain. It verifies both
three-cycles, all $992+1859$ three-cocycle equations, and
the two pairings equal to one, without producer imports
or a matrix-rank calculation. Its complete record is
\path{code/outputs/d_rightzero_six_H3_independently_verified.json}.
The actual quotient from $M$ to either action image is
onto on third homology of the connected graph cofiber:
the natural comparison is $\operatorname{Tor}_1(W,\mathbb F_2)$,
whose two-sided bar has identical degree zero and a
surjection in degree one. Every target one-cycle therefore
lifts to a source one-cycle. This contradicts the third
homology of $\mathbb CP^2$ and excludes both images,
including arbitrary nonfaithful source inflations.

\emph{Three loops and one cross edge.}
Suppose there are loops $a,b$ at vertex zero, one loop
$c$ at vertex one, and the unique cross edge.
There is no swap, and both constant vertex types are
necessary. The constant-one actor forces
$\lambda(a),\lambda(b)\in\{0,\lambda(c)\}$,
with $\lambda(c)\ne0$. The constant-zero actor forces
at least one of the two values to equal $\lambda(c)$.
If one is zero, its loop and its vertex form a right-invariant
collapse: no actor can send the zero-valued loop to a
nonzero-valued loop. Collapsing it gives the connected
five-state graph and an isomorphism on the rank-one
coinvariant space. If both values equal $\lambda(c)$,
identify $a=b$ instead. This is again equivariant; a
constant-one actor kills their difference, so the map
preserves the same coinvariant space. The rational
connected five-state obstruction finishes both cases.

These alternatives exhaust every connected graph with
two vertices and four edges, proving the theorem.
\end{proof}

\begin{remark}[Scope of the image controls]
\label{C:right-zero-six-connected-control-scope}
The two small third-homology certificates have passed their
standalone independent verification. They assert nonzero
third homology of precisely the two displayed literal
categories; the comparison in the proof supplies the
separate nonfaithful-source consequence.
The separate ninety-three-case faithful-image search on a
five-state component plus an isolated core vertex is not used
here and supplies no nonfaithful-source completeness claim.
\end{remark}
